% Source: https://artofproblemsolving.com/community/c3h2134831_mathcounts_chapter_sprint_round_problem_26__2020
% Stable identity: aops-topic-2134831-q01
Place the prism with edges of lengths $5$, $3$ and $4$ along the three coordinate directions, so that the cube faces lie on the planes $x=0,1,\dots,5$, $y=0,1,2,3$ and $z=0,1,\dots,4$.

\begin{asy}
usepackage("amsmath");
size(6cm);
defaultpen(fontsize(10pt));
real X=5, Y=3, Z=4; pair ox=(1,0), oy=(0.42,0.3), oz=(0,1);
pair p(real x,real y,real z){return x*ox+y*oy+z*oz;}
for(int i=0;i<=5;++i){ draw(p(i,0,0)--p(i,0,Z),gray(0.6)); draw(p(i,0,Z)--p(i,Y,Z),gray(0.6)); }
for(int k=0;k<=4;++k){ draw(p(0,0,k)--p(X,0,k),gray(0.6)); draw(p(X,0,k)--p(X,Y,k),gray(0.6)); }
for(int j=0;j<=3;++j){ draw(p(0,j,Z)--p(X,j,Z),gray(0.6)); draw(p(X,j,0)--p(X,j,Z),gray(0.6)); }
draw(p(0,0,0)--p(X,0,0)--p(X,0,Z)--p(0,0,Z)--cycle,linewidth(1));
draw(p(0,0,Z)--p(0,Y,Z)--p(X,Y,Z)--p(X,0,Z),linewidth(1));
draw(p(X,0,0)--p(X,Y,0)--p(X,Y,Z),linewidth(1));
real x0=1,x1=4,y0=0,y1=2,z0=1,z1=3;
fill(p(x0,0,z0)--p(x1,0,z0)--p(x1,0,z1)--p(x0,0,z1)--cycle,palered+opacity(0.6));
pen rp=red+linewidth(1.1);
draw(p(x0,y0,z0)--p(x1,y0,z0)--p(x1,y0,z1)--p(x0,y0,z1)--cycle,rp);
draw(p(x0,y0,z1)--p(x0,y1,z1)--p(x1,y1,z1)--p(x1,y0,z1),rp);
draw(p(x1,y0,z0)--p(x1,y1,z0)--p(x1,y1,z1),rp);
draw(p(x0,y0,z0)--p(x0,y1,z0)--p(x1,y1,z0),red+dashed); draw(p(x0,y1,z0)--p(x0,y1,z1),red+dashed);
label("$5$",p(2.5,0,0),S); label("$3$",p(X,1.5,0),SE); label("$4$",p(0,0,2),W);
for(int i=0;i<=5;++i) dot(p(i,0,0),blue);
label("\mbox{}",(min(currentpicture,user=true)+max(currentpicture,user=true))/2,black); // keeps the text after the figure black when LaTeX typesets the labels inline
\end{asy}

A rectangular prism made of unit cubes is determined by its two faces perpendicular to each direction: a pair of the $6$ planes $x=0,\dots,5$ (the blue dots mark them along the front edge), a pair of the $4$ planes $y=0,\dots,3$, and a pair of the $5$ planes $z=0,\dots,4$. Conversely, every such choice of three pairs gives exactly one prism. (The red prism, for example, uses $x=1,4$, $y=0,2$ and $z=1,3$.) So the number of prisms is
\[\binom62\binom42\binom52=15\cdot6\cdot10=\boxed{900}.\]
This is the three-dimensional version of counting rectangles in a grid by choosing two horizontal and two vertical lines.
